%========================================
% Section：[7] Dhammābhigīti
%========================================
% title-en: Chant in Praise of the Dhamma
\section*{e7 | Dhammābhigīti}

\begin{chantsubsection}
  \chantline
    {Handa mayaṃ dhammābhigītiṃ karoma se}
    {Now let us recite a chant in praise of the Dhamma.}

  \chantline
    {Svākkhātatādiguṇayogavasena seyyo}
    {The Dhamma is supreme, endowed with such qualities as being well expounded,}

  \chantline
    {Yo maggapākapariyattivimokkhabhedo}
    {distinguished as Path, Fruit, Study, and Liberation.}

  \chantline
    {Dhammo kulokapatanā tadadhāridhārī}
    {The Dhamma prevents the fall into the bad world and supports those who uphold it.}

  \chantline
    {Vandāmahaṃ tamaharaṃ varadhammametaṃ}
    {I revere that excellent Dhamma which dispels darkness.}

  \chantlinepair
    {Dhammo yo sabbapāṇīnaṃ}
    {Saraṇaṃ khemamuttamaṃ}
    {The Dhamma is for all beings the supreme, secure refuge.}

  \chantlinepair
    {Dutiyānussatiṭṭhānaṃ}
    {Vandāmi taṃ sirenahaṃ}
    {It is the second foundation of recollection; to it I pay homage with my head.}

  % footnote 1
  \chantlinepair
    {Dhammassāhasmi dāso\footnotemark{} va}
    {Dhammo me sāmikissaro}
    {I am a servant of the Dhamma; the Dhamma is my lord and master.}
  \footnotetext{For ladies: change \textit{dāso} to \textit{dāsī}.}

  \chantlinepair
    {Dhammo dukkhassa ghātā ca}
    {Vidhātā ca hitassa me}
    {The Dhamma is the destroyer of suffering and the establisher of my welfare.}

  \chantlinepair
    {Dhammassāhaṃ niyyādemi}
    {Sarīrañjīvitañcidaṃ}
    {To the Dhamma I entrust this body and this life.}

  % footnote 2
  \chantlinepair
    {Vandantohaṃ\footnotemark{} carissāmi}
    {Dhammasseva sudhammataṃ}
    {Paying homage, I will live devoted only to the true nature of the Dhamma.}
  \footnotetext{For ladies: change \textit{Vandantohaṃ} to \textit{Vandantīhaṃ}.}

  \chantlinepair
    {Natthi me saraṇaṃ aññaṃ}
    {Dhammo me saraṇaṃ varaṃ}
    {I have no other refuge; the Dhamma is my excellent refuge.}

  \chantlinepair
    {Etena saccavajjena}
    {Vaḍḍheyyaṃ satthu sāsane}
    {By this truthful utterance, may I grow in the Teacher's dispensation.}

  % footnote 3
  \chantlinepair
    {Dhammaṃ me vandamānena\footnotemark{}}
    {Yaṃ puññaṃ pasutaṃ idha}
    {Whatever merit I have gathered here by revering the Dhamma,}
  \footnotetext{For ladies: change \textit{vandamānena} to \textit{vandamānāya}.}

  \chantlinepair
    {Sabbepi antarāyā me}
    {Māhesuṃ tassa tejasā}
    {may all my obstacles be dispelled through its power.}
\end{chantsubsection}
